\chapter{Related work}
\label{chap:related-work}

This chapter reviews the four bodies of work this thesis draws on: benchmarks of language-model
inference on \glspl{sbc}, post-training quantisation, grammar-constrained decoding, and
spoken-language understanding. Each section states what its literature establishes and what it
leaves unmeasured for this task and this hardware. Section~\ref{sec:positioning} combines the four
into the gap that RQ1 addresses and maps each contribution onto it.

\section{Edge language-model inference and single-board-computer benchmarking}
Running a language model without a network connection or a data-centre accelerator is now an
active subject of measurement. A recent benchmark evaluates twenty-five quantised language
models on three \glspl{sbc} under two inference runtimes, Ollama and
Llamafile~\cite{sbc2025}. It finds that such boards reliably support models of up to roughly
1.5~billion parameters, with larger models dropping to a few tokens per second. On the
one board where both runtimes were compared, Llamafile reached up to four times the throughput of
Ollama~\cite{sbc2025}. That result sets the shape of the problem this thesis inherits. On
\gls{cpu}-only edge hardware throughput is the scarce resource, and the models that fit within it
are small ones; the inference runtime is a lever on it of the same order as the choice of model.

What that benchmark does not report is what happens to a model once it is asked to do something
with the tokens it produces. Its outcome metrics describe speed and resource use, measured
independently of any downstream task. A configuration that clears a throughput floor and one that
clears it while producing ungrammatical or semantically wrong output are therefore
indistinguishable in its tables. This thesis narrows the scope to one task, one device, and four
fine-tuned models (Qwen2.5-0.5B~\cite{qwen25}, SmolLM2-360M~\cite{smollm2},
Llama-3.2-1B~\cite{llama32}, and H2O-Danube3-500M~\cite{danube3}). In exchange it measures what
the benchmark does not: \gls{em}, intent and slot F1, schema validity, and false-command rate.
Three of those models are carried through to quantised artefacts and timed against the
20~tok/s throughput floor of Chapter~\ref{chap:introduction}. The fourth, a parameter-matched
control, is reported at fp16 only, because llama.cpp segments its prompt differently from the
tokeniser it was trained under (Chapter~\ref{chap:method}). The accuracy cost of meeting a
throughput budget on an \gls{sbc} is thus treated as a quantity to be measured, not
assumed from a benchmark that reports throughput alone.

The runtime that makes this measurement possible on the target hardware is llama.cpp, a C/C\,++
inference engine for quantised language models that runs on a \gls{cpu} without external
dependencies, and in particular without a \gls{gpu} driver~\cite{llamacpp}. Its \gls{gguf}
weight format and its quantisation schemes shrink the weights a Raspberry~Pi~5 must read for every
decoded token; Section~\ref{sec:quantisation} takes up what those schemes cost in accuracy.

\section{Quantisation}
\label{sec:quantisation}
% NOTE: \cite{llamacpp} below is correct and load-bearing, not a placeholder. No
%       peer-reviewed analysis of the k-quant scheme itself appears to exist, so the
%       runtime is the primary source for the SCHEME DESCRIPTION. \cite{gptq} and
%       \cite{awq} carry the field-level claims around it. Do not 'fix' this.
Post-training quantisation is the established route to reducing a trained model's memory
footprint~\cite{gptq}, and on-device deployment is one of its stated motivations~\cite{awq}.
This thesis deploys two of llama.cpp's \gls{gguf} quantisation formats. Q8\_0, a legacy format,
stores weights in blocks of 32 at eight bits with one scale per block, 8.5~bits per weight in
all. Q4\_K\_M belongs to the K-quant family, which stores weights in super-blocks of 256 with
quantised per-sub-block scales, 4.5~bits per weight at the four-bit level~\cite{llamacpp}. Its
target is four bits for most tensors. A fixed rule raises the output projection, and a subset of
the attention-value and feed-forward down-projection tensors chosen by layer position, to six
bits~\cite{llamacpp}. A K-quant also requires each tensor's rows to divide into whole
super-blocks: a tensor whose row length is not a multiple of 256 falls back to a legacy format,
five bits in place of four and eight in place of six~\cite{llamacpp}. The name Q4\_K\_M therefore
denotes an allocation recipe rather than a bit width, and in a model whose hidden dimension is not
a multiple of 256 most of the weights end up at five or eight bits. The recipe allocates bits by
tensor role, position and shape, not by any measurement of error. The method literature measures
instead: it finds that protecting the roughly one percent of weight channels that activation
statistics mark as salient greatly reduces quantisation error~\cite{awq}.

The runtime's documentation states that the accuracy cost of its formats is usually measured as
a perplexity or a Kullback--Leibler divergence against the unquantised model~\cite{llamacpp},
not as a change in accuracy on a downstream task. A perplexity number averaged over held-out text
says how much less confident the quantised model is about the next token in general. It does not
say whether the intent assigned to a spoken command, or the slot values filled in, still match the
ground truth once the weights have been rounded to lower precision. This thesis needs the
second number, because its model's output is acted on by a swarm controller after automatic
validation, not read by a person. The papers introducing post-training quantisation methods report
their cost mainly in generic terms: SpQR reports a relative perplexity loss of under one
percent~\cite{spqr}, and GPTQ reports perplexity together with zero-shot accuracy on generic
benchmarks such as LAMBADA, ARC and PIQA~\cite{gptq}.

Evaluation studies have since measured downstream accuracy under quantisation at
scale~\cite{kurtic2025}, and for llama.cpp's own formats specifically~\cite{kurt2026}. Their
models are mostly of eight billion parameters and above, scored on general-purpose reasoning,
knowledge and coding suites with the decoder unconstrained; the smallest model either study
evaluates has 1.5~billion parameters~\cite{kurtic2025}. Small models have also been studied
directly: a benchmark of quantisation at sub-billion scale, which includes Qwen2.5-0.5B, finds that
small models differ from large ones in their sensitivity to quantisation, so that methods tuned on
large models transfer poorly~\cite{slmquant}. None of these studies scores a structured
command-parsing task by \gls{em}, intent and slot F1, or schema validity, and none decodes under
a grammar. Chapter~\ref{chap:results} measures the accuracy cost of the whole deployment pipeline
on that task, for the three quantised models at both levels, as part of Contribution~C3.
Section~\ref{sec:quantisation-procedure} explains why that pipeline cost is not the cost of
quantisation alone, and why it is not to be read against the studies above.

\section{Constrained decoding}
\label{sec:constrained-decoding}
Grammar-constrained generation restricts, at every decoding step, the set of tokens the model is
allowed to sample to those consistent with a formal grammar, so that no output can leave the
grammar's language regardless of what the model would otherwise have produced. The technique has
been applied to structured language-processing tasks without fine-tuning~\cite{geng2023}. Its
formulations compile a regular expression or a context-free grammar into an automaton over the
vocabulary and mask the decoder against it~\cite{willard2023,koo2024}, with correctness results
for the language classes such a construction covers~\cite{koo2024}. \Gls{gbnf} is the grammar
formalism llama.cpp implements it with: a context-free grammar applied at each decoding step to
mask the tokens it cannot accept. Its per-step cost depends on the grammar's
structure, and the documentation warns that some rule shapes make sampling slow~\cite{llamacpp}.
This is a different mechanism from prompting a model to produce well-formed output and then
checking the result: the grammar acts inside the sampling loop, before an invalid token can be
emitted, not after a full sequence has been generated.

That difference in mechanism is also a difference in the kind of claim that can be made about
structural validity. Without a grammar, whether a model's output parses against a schema is an
empirical question, answered by counting how often it does so on some test set, and the answer
can change whenever the input distribution changes. With a grammar compiled from the same schema,
the question is closed by construction: no string outside the grammar's language is reachable, on
any input, provided the generation cap exceeds the longest string that language contains. The
bounded grammar of Section~\ref{sec:schema} makes that length finite. Constrained decoding thus
moves structural validity from a rate that fine-tuning and a favourable test set can inflate to a
property the decoder enforces independently of both. The grammar ablation of
Chapter~\ref{chap:results} decodes the same fine-tuned models with and without the grammar, and so
measures how often they leave the format when nothing prevents them.

That property has a cost, which is why the ablation reports accuracy as well as validity.
Constraining the decoder alters the distribution the model samples from: outputs remain
grammatical, but their relative likelihoods no longer track what the unconstrained model would
have assigned, so a grammar can make a structurally valid command more likely without making the
correct one more likely~\cite{park2024}. Format restriction has separately been observed to
degrade accuracy on reasoning tasks, more so as the restriction tightens, while classification
tasks were in some cases helped by it~\cite{tam2024}. The size of either effect on a given task
does not follow from that literature. The distortion result establishes that the effect exists,
not how large it is, and the format-restriction study finds even its direction task-dependent.
Of the two task types it compares, classification is the closer to parsing a command into one of
ten intents. Chapter~\ref{chap:results} therefore reports accuracy with and without the grammar,
not structural validity alone.

\section{Spoken-language understanding for robotics}
\label{sec:slu-robotics}
Parsing an utterance into an intent and a set of slot values is an established formulation in
spoken-language understanding, not one this thesis introduces. A survey of the field treats the
extraction of this semantic frame as the object of the task, and classifies methods by whether they
model intent and slots separately or jointly~\cite{qin2021}. MASSIVE, a dataset of one million
utterances in 51 languages labelled with 60 intents and 55 slot types, is built around the same
formulation~\cite{massive}. Its baselines, encoders of 258 to 580~million parameters, reach
locale-averaged intent accuracies of 85.1--86.1\% and slot F1 of 73.6--76.8\%, but \gls{em}
of only 63.7--66.6\%~\cite{massive}. For the single-utterance commands considered here, a spoken
instruction to a robot has the same shape: one intent from a fixed set and a handful of slot
values. Language-model control of robots and aircraft more broadly is reviewed in the
\emph{M\'emoire d'Ing\'enieur}.

What differs in this deployment is the setting, not the shape of the task. MASSIVE's baselines are
trained on data-centre \glspl{gpu}, decoded without a grammar, and scored with no latency or
memory limit~\cite{massive}. Their lowest figure is \gls{em}, and it is the metric a controller
depends on, since a command with one wrong slot is a wrong command. Here the model has to
run offline, at sub-billion-parameter scale, within the throughput budget of a Raspberry~Pi~5, and
its output is acted on by the swarm controller after automatic validation, with no person reading
it first. No existing corpus covers this command vocabulary (Section~\ref{sec:dataset}). That the
task shape is well studied does not show that it can be solved at this size and within this
budget, with an error rate low enough for output that no person reads. The evaluation protocol of
Chapter~\ref{chap:method} is built to answer that question.

\section{Positioning}
\label{sec:positioning}

\paragraph{Synthesis.} Each of the four literatures reviewed in this chapter answers its own
question. Benchmarking on \glspl{sbc} has found that such boards reliably serve models of up to
about 1.5~billion parameters, and that on one board the choice of runtime changed throughput by up
to four times~\cite{sbc2025}. Post-training quantisation is the established way to shrink a trained
model~\cite{gptq}, with on-device deployment among its motivations~\cite{awq}. Its cost is reported as perplexity and generic zero-shot accuracy by
the papers introducing the methods~\cite{spqr,gptq}, and as accuracy on general-purpose suites by
the studies evaluating them, including at sub-billion scale~\cite{kurtic2025,kurt2026,slmquant}.
Grammar-constrained decoding makes a model's output structurally valid by construction rather than
by the accuracy of the model that produced it~\cite{willard2023,koo2024}, at costs its own
literature documents: a distorted sampling distribution~\cite{park2024} and, on reasoning tasks,
lower accuracy~\cite{tam2024}. Intent and slot parsing
is an established task formulation, with a large multilingual benchmark built for
it~\cite{massive}. Some pairs of these literatures have already been combined: quantisation with
\gls{cpu} inference and downstream accuracy~\cite{kurt2026}, quantisation with small
models~\cite{slmquant}, and format restriction with task accuracy~\cite{tam2024}. No study
reviewed here combines all four.

\paragraph{Gap.} The gap lies in that combination: a sub-billion model fine-tuned to a command
schema, quantised to a format an \gls{sbc} can serve, decoded under a grammar, and
scored on task metrics alongside its throughput on the board. The single-board-computer benchmark
reviewed here reports throughput without reference to any task, so it cannot say what a
configuration that meets a throughput floor costs in task accuracy. The quantisation studies report
accuracy on general-purpose benchmarks with the decoder unconstrained, not \gls{em}, intent and
slot F1, or schema validity on a command-parsing task. The constrained-decoding studies reviewed here
evaluate grammars and format restrictions on models that were not fine-tuned to the target format, so what a grammar still adds
once fine-tuning has taught a model that format is not reported. Intent-and-slot benchmarks such as
MASSIVE are scored on data-centre hardware without a latency budget or a grammar, and no existing
corpus covers this command vocabulary. Choosing a model to deploy on this hardware requires all
four quantities measured together, on the configuration that is actually deployed.

\paragraph{Delta.} The three contributions of this thesis address that gap as follows.
Contribution C1, the frozen command schema and its matched, bounded \gls{gbnf} grammar, makes
structural validity on this task a property of the decoder. The grammar ablation, run with the
harness of Contribution C3, decodes the same fine-tuned, quantised models with and without that
grammar, and so reports what the grammar adds once fine-tuning has taught the models the format.
Contribution C2, the label-first dataset (each label written before the sentences that express
it), supplies what a general-purpose corpus such as MASSIVE does not: training and test data for
this command vocabulary, split at the template-family level. Contribution C3, the
quantised-inference harness, measures \gls{em}, intent and slot F1, schema validity, and
false-command rate for three of the four fine-tuned models at Q8\_0 and Q4\_K\_M, with accuracy
scored under the deployed runtime and grammar and throughput timed on the Raspberry~Pi~5 at
Q4\_K\_M. It reports the accuracy cost of the whole deployment pipeline on this task at
sub-billion scale, next to the throughput that single-board-computer benchmarking measures alone.
Together they supply the measurements from which Chapter~\ref{chap:results} answers RQ1 for this
task and this device. Chapter~\ref{chap:method} specifies how
each contribution was built, and Chapter~\ref{chap:results} reports what it measured.
